\chapter{Pendahuluan}

\section{Latar Belakang}
Pembakaran adalah suatu proses kompleks yang memanaskan dan mengoksidasi bahan bakar dengan cepat, menghasilkan panas pada prosesnya. Api merupakan hasil pembakaran yang berkelanjutan sendiri yang ditandai dengan emisi panas dan disertai dengan nyala api dan/atau asap \cite{scott2012wildfire}. Kebakaran terbentuk dari fragmen-fragment mudah terbakar yang terbentuk dari sumber api asli. Fragmen-fragmen tersebut membuat kebakaran dapat menyebar dengan cepat \cite{10.3389/fmech.2022.819633}.

Analisis DPR RI (2025) menekankan bahwa Peningkatan fungsi perkotaan dan perubahan iklim membuat permukiman padat penduduk di perkotaan Indonesia sangat rentan terhadap ancaman kebakaran. Fungsi ekologi kota melemah sehingga kebakaran lebih sulit dikendalikan dan dampaknya bisa signifikan. Untuk mencegah dampak bencana kebakaran menjadi terlalu parah, tindakan-tindakan preventif diperlukan \cite{prihatin2025mitigasi}.

\textit{Fire drill} yang dilaksanakan secara efektif dapat mengurangi kebingungan, mengurangi cedera, dan menyelamatkan nyawa dengan mengembangkan pemahaman mengenai rute-rute keluar, prosedur keselamatan dan perilaku kelompok yang terkoordinasi. Di lingkungan dengan kepadatan tinggi, kesiapan dapat mempengaruhi tingkat keselamatan dan efisiensi tanggapan darurat secara signifikan\cite{michael2024measuring}. 

Selain \textit{fire drill}, metode-metode yang umum digunakan untuk pelatihan evakuasi kebakaran adalah presentasi, video pembelajaran, TV, dan media fisik.  Pengalaman langsung umumnya sulit dilakukan karena alasan keamanan. Namun, di lingkungan dengan kebakaran yang asli, psikologi dan perilaku manusia akan sangat berpengaruh sehingga pelatihan umum yang bersifat konvensional tidak terlalu efektif. Tidak hanya \textit{fire drill}, pelatihan APAR (alat pemadam api ringan) juga salah satu upaya penting. Sebuah insiden api laboratorium di Northwestern University menyorot fakta bahwa para asisten yang tidak terlatih menggunakan pemadam api menunda respon aktif \cite{doi:10.1021/acs.chas.2c00094}. Namun, dalam praktiknya, pelatihan yang menerapkan lingkungan kebakaran yang nyata membutuhkan biaya yang terlalu tinggi, tidak dapat digunakan ulang, dan juga sulit dikendalikan\cite{9824639}. 


Teknologi Virtual Reality dapat digunakan dalam pembelajaran evakuasi sebagai alternatif.  Selain imersif, teknologi ini lebih efisien secara biaya karena pelatihan dapat dilakukan kapan saja. Risiko finansial karena kerusakan akibat \textit{fire drill} dan pelatihan secara fisik juga dapat dicegah karena VR menggunakan aset virtual yang tidak merusak\cite{9256834}.
Karena keunggulan tersebut, teknologi VR telah digunakan secara ekstensif di berbagai domain profesional dan juga memiliki target populasi yang beragam. Prosedur medis kompleks membutuhkan presisi dan keamanan prosedural yang tinggi. Para ahli bedah melatih kemampuan melakukan prosedur kompleks seperti \textit{laparoscopic} atau \textit{robot-assisted surgery} dengan bantuan VR. Pelatihan keamanan dengan VR juga dilakukan oleh para pekerja konstruksi untuk melatih kemampuan identifikasi bahaya dan operasi alat berat secara aman\cite{Stefan2023}. 

\section{Rumusan Masalah}

Fasilitas pendidikan seperti gedung Smart Green Learning Center (SGLC) UGM juga memiliki risiko kebakaran sehingga diperlukan upaya preventif untuk mencegah adanya korban dan signifikansi kerugian. Oleh karena itu, kemampuan sivitas akademika (dosen, mahasiswa, dan tenaga kependidikan) dapat ditingkatkan dengan menggunakan simulasi kebakaran berbasis VR. Masalah dapat dirumuskan sebagai berikut:
\begin{enumerate}
	\item Sebuah metode pembelajaran kebencanaan (kebakaran) dengan menggunakan VR perlu dikembangkan
	\item Perbandingan antara metode VR dan metode yang sudah ada perlu dilakukan
\end{enumerate}


perbandingan dengan metode usability yesting


%-------------------------------------------------	

\section{Tujuan Penelitian}

% Tujuan penelitian pada skripsi Teknik (TE, TB, TIF) adalah menentukan sasaran atau target yang ingin dicapai melalui proses penelitian. Tujuan penelitian bisa beragam sesuai dengan bidang ilmu yang dipelajari, topik penelitian, dan permasalahan yang akan dicari solusinya.

%Secara umum, tujuan penelitian pada skripsi bidang teknik adalah untuk:
%
%\begin{itemize}
%	\item Mengidentifikasi dan menganalisis masalah atau permasalahan dalam bidang \textit{engineering}.
%	\item Meningkatkan pemahaman dan wawasan tentang bidang \textit{engineering} melalui penerapan teori dan metodologi yang sesuai.
%	\item Mengembangkan solusi atau produk baru yang inovatif dan efektif dalam 
%	bidang \textit{engineering}.
%	\item Menunjukkan penerapan prinsip-prinsip keteknikan dalam solusi atau produk yang dikembangkan.
%	\item Menjelaskan implikasi dan rekomendasi dari hasil penelitian bagi bidang \textit{engineering} dan masyarakat.
%\end{itemize}
%
%
%\noindent Catatan: Tujuan penelitian dalam skripsi bidang teknik harus jelas, spesifik, terukur, dan dapat dicapai dalam jangka waktu yang ditentukan.


Tujuan dari pengembangan tugas akhir ini adalah:
\begin{enumerate}
	\item Mengembangkan aplikasi VR yang mensimulasikan penggunaan APAR di gedung SGLC UGM
	\item Menguji fungsionalitas, kebergunaan, dan pengalaman pengguna dari media pembelajaran elektronik terhadap calon pengguna
	\item Membandingkan capaian pembelajaran metode (pilih satu metode konvensional) dan aplikasi pembelajaran dengan VR
\end{enumerate}


%-------------------------------------------------	


\section{Batasan Penelitian}

Berikut adalah batasan penelitian dari tugas akhir ini:

\begin{enumerate}
	\item Objek Penelitian : Penelitian berfokus pada XXX dan XXX dalam pelatihan evakuasi kebakaran menggunakan teknologi *Virtual Reality* (VR) di gedung SGLC Fakultas Teknik Universitas Gadjah Mada. Ini bertujuan untuk XXX
	\item 	Metode Penelitian : Penelitian ini menggunakan metode XXX dengan aplikasi VR. Setelah mencoba menggunakan VR, peserta akan XXX. Hasilnya akan XXX
	\item Waktu dan Tempat Penelitian : Penelitian ini akan dilaksanakan selama 6 bulan di gedung SGLC Fakultas Teknik UGM
	\item Populasi dan Sampel : Populasi penelitian ini adalah mahasiswa UGM. Sampel yang akan diteliti terdiri dari XX partisipan yang dipilih secara acak dari populasi tersebut untuk melakukan pelatihan evakuasi menggunakan aplikasi VR
	\item Keterbatasan Penelitian :
	\begin{enumerate}
		\item Sumber data yang terbatas : Aplikasi yang dikembangkan hanya dapat digunakan di Meta Quest 2 karena keterbatasan dalam peralatan VR yang tersedia
		\item Waktu : Waktu untuk pelaksanaan, penelitian, dan pengumpulan data cukup terbatas
		\item Metodologi : Generalisasi hasil cukup terbataskarena penelitian hanya dilakukan di satu gedung dengan populasi tertentu
		\item Kemampuan Peneliti : Penliti memiliki keterbatasan dalam penglaman dan kemampuan mengelola dan menganalisisi data simulaasi VR.n aplikasi VR yang mensimulasikan evakuasi kebakaran di gedung SGLC UGM

	\end{enumerate}
   
	
\end{enumerate}

\section{Manfaat Penelitian}

Pengembangan aplikasi VR diharapkan dapat meningkatkan kemampuan sivitas akademika Fakultas Teknik UGM dalam melakukan evakuasi kebakaran. Selain itu, penggunaan VR diharapkan dapat meningkatkan efektivitas pembelajaran evakuasi bencana di Fakultas Teknik UGM.


\section{Sistematika Penulisan}

Bab I menjelaskan latar belakang, rumusan masalah yang akan dijawab dalam penelitian, batasan permasalahan dalam pelaksanaan penelitian, tujuan yang akan dicapai dalam penelitian ini, dan manfaat penelitian bagi pihak-pihak terkait.

Bab II memuat hasil tinjauan literatur yang mendukung penelitian berdasarkan penerapan VR dalam pelatihan dan pembelajaran yang menjadi dasar teori pendukung dalam pengembangan aplikasi VR. Bagian ini juga akan menjelaskan teori dasar dan terminologi yang perlu dipahami untuk mengembangkan aplikasi dan memahami cara kerja aplikasi.

Bab III menjabarkan persyaratan yang diperlukan penulis dalam mengembangkan aplikasi VR seperti penjelasan rinci perangkat lunak dan keras yang digunakan, urutan langkah pengembangan aplikasi, serta metode pengujian.

Bab IV menjelaskan keluaran pengembangan aplikasi dan proses pengembangan mulai dari desain, sistem sistem final, dan hasil pengujian [akan disajikan dalam bentuk gambar] yang dilengkapi deskripsi penjelasan tiap tahapannya

Bab V memuat rangkuman seluruh proses penerapan VR dalam pelatihan evakuasi bencana kebakaran dan keberhasilan aplikasi dalam menyelesaikan masalah yang telah dirumuskan. Bagian ini juga memuat rekomendasi untuk mengoptimalkan keluaran yang dapat menjadi sumber pembelajaran dalam pengembangan sistem pelatihan kebakaran di masa depan
